\section{Fundamentos de las proteínas: secuencia, estructura y función}

\subsection{Las proteínas son los "actuadores" de la biología}

Las proteínas (protein) son una clase de macromoléculas biológicas extremadamente diversas, que pueden considerarse \textbf{los actuadores de la biología} (biology's actuators): impulsan y regulan casi todas las funciones dentro de la célula. El rango de funciones de las proteínas es sumamente amplio e incluye, entre otras:

\begin{itemize}
    \item \textbf{Enzimas} (enzyme): catalizan reacciones químicas específicas, incluida la edición a nivel del genoma (como las proteínas asociadas a CRISPR)
    \item \textbf{Proteínas estructurales}: constituyen el citoesqueleto, la matriz extracelular y otras estructuras
    \item \textbf{Proteínas de señalización}: transmiten señales entre células y dentro de ellas (como las proteínas que se unen al calcio)
    \item \textbf{Proteínas terapéuticas}: como las proteínas farmacológicas ya aprobadas clínicamente para tratar la leucemia
    \item \textbf{Proteínas de transporte}: transportan moléculas dentro y fuera de la célula
\end{itemize}

\begin{knowledgebox}{Las proteínas como fármacos terapéuticos}
Las proteínas no solo son ejecutoras naturales de funciones biológicas, también pueden \textbf{diseñarse mediante ingeniería} como fármacos terapéuticos. Por ejemplo, los medicamentos con anticuerpos (un tipo especial de proteína) se usan ampliamente en inmunoterapia contra el cáncer. Amini mostró en la clase un fármaco proteico ya aprobado clínicamente para tratar un tipo específico de leucemia, lo que demuestra que el valor de traslación clínica del diseño de proteínas ya se ha verificado.
\end{knowledgebox}

\subsection{Los tres niveles de representación de las proteínas}

La biología de las proteínas puede entenderse mediante tres niveles, que forman una relación jerárquica desde la codificación lineal hasta la estructura tridimensional y, finalmente, la función biológica:

\begin{table}[H]
\centering
\caption{Los tres niveles de representación de las proteínas}
\begin{tabular}{@{}lll@{}}
\toprule
\textbf{Nivel} & \textbf{Descripción} & \textbf{Características de los datos} \\
\midrule
Secuencia (Sequence) & Disposición lineal de aminoácidos & Secuencia de símbolos discretos, alfabeto de 20 aminoácidos \\
Estructura (Structure) & Conformación tridimensional en el espacio & Coordenadas continuas, información geométrica a nivel atómico \\
Función (Function) & Actividad biológica y efecto & Anotaciones diversas, difíciles de cuantificar \\
\bottomrule
\end{tabular}
\end{table}

\textbf{Nivel de secuencia}: cada proteína queda definida por una secuencia de aminoácidos (amino acid). En la naturaleza existen 20 aminoácidos estándar, cada uno representable con un código de una sola letra (por ejemplo, A para alanina, Alanine, y G para glicina, Glycine). Por ello, la secuencia de una proteína puede considerarse un "lenguaje" basado en un vocabulario de 20 caracteres, lo que ofrece una analogía natural para aplicar técnicas de PLN al procesamiento de datos de proteínas.

\textbf{Nivel de estructura}: la secuencia de aminoácidos se pliega (folding) mediante procesos fisicoquímicos complejos hasta adoptar una estructura tridimensional específica. La estructura determina cómo interactúa la proteína con otras moléculas y constituye la base física para la realización de su función.

\textbf{Nivel de función}: la función biológica de una proteína (como la catálisis, la unión o la transducción de señales) suele estar determinada por su estructura, y esta a su vez está codificada por la secuencia.

\begin{importantbox}{Ingeniería inversa del dogma central}
La investigación tradicional sobre proteínas sigue la vía directa de secuencia $\rightarrow$ estructura $\rightarrow$ función. Sin embargo, el diseño de proteínas impulsado por IA busca \textbf{invertir este paradigma}: partir de la \textbf{función} que deseamos obtener y deducir de manera inversa la \textbf{secuencia} que la realiza. Este es precisamente el reto central del diseño generativo de proteínas:
$$\text{Function} \rightarrow \text{Sequence} \rightarrow \text{Structure (验证)}$$
\end{importantbox}

\subsection{Herramientas de IA existentes para proteínas}

En el campo de la IA aplicada a proteínas ya existe una cantidad considerable de trabajo dirigido a cada uno de los eslabones del pipeline secuencia-estructura-función:

\begin{itemize}
    \item \textbf{Modelos de lenguaje de proteínas} (Protein Language Models): como la serie ESM, que aprenden representaciones de secuencias de proteínas y capturan información evolutiva y funcional
    \item \textbf{Modelos generativos de estructuras}: como RFdiffusion, capaces de diseñar en el espacio estructural nuevas geometrías que no existen en la naturaleza
    \item \textbf{Predicción de secuencia a estructura}: como \textbf{AlphaFold}, que dada una secuencia de aminoácidos predice su estructura tridimensional; este es uno de los avances más importantes de los últimos años en biología computacional
\end{itemize}

\begin{warningbox}{La brecha de escala entre los datos de estructura y los datos de secuencia}
En el campo de las proteínas existe una \textbf{asimetría de datos} considerable:
\begin{itemize}
    \item \textbf{Datos de secuencia}: cientos de millones, o incluso miles de millones, de secuencias de proteínas distintas (provenientes de la secuenciación de genomas)
    \item \textbf{Datos de estructura}: apenas unas 300,000 estructuras de proteínas resueltas experimentalmente (provenientes de la base de datos PDB)
\end{itemize}
Aunque AlphaFold puede predecir estructuras, se trata de una \textbf{predicción computacional} y no de una medición experimental. Las proteínas en las bases de datos de estructuras tienden a ser globulares (globular) y solubles, con abundancia de elementos estructurales de tipo $\alpha$-hélice, lo que introduce un sesgo sistemático en los modelos entrenados con datos de estructura.
\end{warningbox}

\subsection{Resumen de la sección}

Las proteínas son la macromolécula biológica central que conecta la codificación de secuencias con la función biológica. Comprender la relación jerárquica de tres niveles entre secuencia, estructura y función es la base para construir herramientas de IA de diseño de proteínas. Las herramientas existentes ya son capaces de abordar cada eslabón del pipeline por separado, pero lograr el diseño inverso de función a secuencia sigue siendo un reto abierto. En cuanto a los datos, la riqueza de los datos de secuencia supera con creces a la de los datos de estructura, y esta asimetría influye profundamente en la elección de las estrategias de diseño de los modelos.
